\chapter{분포 비교: 평균, 분산, 분포 전체를 검정으로 가르기}
\label{ch:compare-dist}
% source: deck 11 slide 1

\begin{inbrief}
  \item[Why] "두 집단이 다르다"는 말에는 세 가지 뜻이 있다. 평균이 다르다,
             퍼짐이 다르다, 모양 전체가 다르다. 뜻마다 맞는 검정이 따로
             있다.
  \item[When] 두 표본(또는 셋 이상의 집단)을 놓고 "같은 분포에서 나왔는가"를
              판정할 때. 데이터 분석에서 가장 흔한 질문이다(슬라이드 2).
  \item[Where] 반바지 대 긴바지 달리기 기록, 두 밭의 토마토 키, 학년별
               만화 시청 시간, 표본이 $N(0,1)$에서 왔는가, HW2의 모형 적합.
  \item[How] 평균이면 \term{t검정}(t-test), 셋 이상 집단의 평균이면
             \term{분산분석}(Analysis of Variance, ANOVA)의
             \term{F검정}(F-test), 분포 전체면
             \term{콜모고로프-스미르노프 검정}(Kolmogorov-Smirnov Test,
             KS 검정) $\to$ 통계량을 계산한다 $\to$ 귀무가설 아래 그
             통계량의 분포에서 p값을 읽는다(\Cref{ch:pvalue}).
\end{inbrief}

이 장의 출처는 \texttt{notes/11\_Comparing Distributions.pptx}(24 슬라이드)
하나다. 그림 속 예제(교사 경력, 토마토, 만화 시청 시간, $N(0,1)$ 100개,
표본 $X$와 $Y$)는 슬라이드 자체를 출처로 적는다. t표, F표, KS 임계값표는
그대로 옮기지 않고 필요한 한 칸만 인용한다.

\analogy{지도 비유}를 이어 쓴다. KS 검정은 "두 지도를 겹쳤을 때 가장 크게
어긋나는 한 지점의 거리"로 요약한다.

% =====================================================================
\section{배경}
\label{sec:cd-background}
% source: deck 11 slide 2

이 장은 다음 개념을 그대로 가져와 쓴다.

\begin{itemize}
  \item 가설검정 절차와 단측/양측 형태(\Cref{fig:ht-flow},
        \Cref{sec:ht-forms}, \Cref{fig:ht-tails}).
  \item p값, 연속 데이터의 p값 = 꼬리 넓이, $z$통계량
        (\Cref{def:pvalue}, \Cref{sec:pv-continuous},
        \Cref{sec:pv-sim-vs-param}). \term{t통계량}(t-statistic)은
        "표준오차를 표본에서 추정한 $z$"라고 보면 된다.
  \item 표준편차, 분산, 제곱거리 합(\Cref{sec:spread}, \Cref{def:spread}).
        ANOVA의 \term{제곱합}(Sum of Squares, SS)이 바로 이것이다.
  \item 표준오차(\Cref{pit:sd-se-ci}): $s/\sqrt n$.
  \item 합친 평균 대 집단 평균(\Cref{sec:ht-two-groups},
        \Cref{ex:ht-distance}): 덱 9가 남긴 숙제, "이 거리를 정식 검정
        통계량으로 만드는 방법"을 이 장이 푼다.
  \item 누적분포함수(\Cref{sec:cdf})를 이미 알고 있다고 본다.
        \term{경험적 누적분포함수}(Empirical Cumulative Distribution
        Function, 경험적 CDF)는 이 장 \Cref{sec:cd-ks}에서 직접 정의한다.
\end{itemize}

새 개념 하나는 이 장에서 직접 정의한다. \term{자유도}(Degrees of
Freedom, df 또는 $\nu$)는 "표본에서 값을 몇 개나 자유롭게 움직일 수
있는가"를 센 것이다. 평균을 추정하는 데 값 하나를 이미 썼으면, 나머지
값 중 하나는 평균에 맞춰 자동으로 정해지므로 자유롭게 움직일 수 있는
값이 하나 줄어든다. 표본 크기 $n$에서 평균 하나를 추정하면 자유도는
$n-1$이다.

\Cref{tab:cd-symbols}가 이 장에서 쓰는 기호를 모은다.

\begin{table}[ht]\centering
\begin{tabularx}{\linewidth}{@{}l X@{}}
\toprule
기호 & 뜻 \\
\midrule
$\bar x$, $s$ & 표본 평균, 표본 표준편차(Sample Standard Deviation) \\
$\mu_0$ & 귀무가설이 정하는 기준값 \\
$t$ & t통계량 \\
$\nu$ 또는 df & 자유도 \\
$m$ & 집단 수 \\
$n$ & 집단당 표본 크기 \\
SSB, SSW & 집단 간(Between), 집단 내(Within) 제곱합 \\
$F$ & F통계량(분산분석용) \\
$F_{\text{obs}}$, $F_{\text{exp}}$ & 관측/기대 누적분포함수(KS용) \\
$D_n$ & KS 통계량, 두 누적분포함수 사이 최대 세로 격차 \\
\bottomrule
\end{tabularx}
\caption{이 장에서 쓰는 기호.}
\label{tab:cd-symbols}
\end{table}

\begin{pitfall}[$F$는 이 장에서 두 가지 뜻으로 쓰인다]
분산분석의 \textbf{F통계량}(위 표의 $F$)과 누적분포함수인
\textbf{$F_{\text{obs}}$, $F_{\text{exp}}$}는 이름만 같을 뿐 다른
대상이다. 아래 첨자가 있으면 누적분포함수, 없으면 분산분석의 통계량으로
읽는다. 슬라이드는 이 둘을 구분 없이 $F$로 적는다.
\end{pitfall}

슬라이드는 표준편차 기호로 $\sigma_x$(모집단 값처럼 보이는 표기)를 쓰지만,
실제 계산은 표본에서 잰 값이다. 이 장은 그 값을 표본 표준편차 $s$로
통일해 적는다.

% =====================================================================
\section{왜 필요한가}
\label{sec:cd-why}
% source: deck 11 slide 2

\begin{itemize}
  \item 문제 1: \Cref{ex:ht-distance}는 "제곱거리가 $0.7\%$밖에 줄지
        않았다"에서 멈췄다. 얼마나 줄어야 "크다"인지 기준이 없었다.
  \item 문제 2: \Cref{subsec:ci-compare}의 구간 비교는 한 방향으로만
        강하다. "겹친다고 같은 것은 아니다"는 말할 수 있어도, 겹칠 때
        판정할 도구가 없었다.
  \item 문제 3: HW2(\Cref{sec:hw2-compare-how})는 $D$를 계산했지만
        "얼마면 기각인가"를 정하지 않았다.
\end{itemize}

처방은 하나다. 차이를 \textbf{표준오차 단위로 잰 통계량}을 만들고,
귀무가설 아래 그 통계량의 분포에서 p값을 읽는다. 무엇을 비교하느냐에
따라 통계량이 셋으로 갈린다. 평균이면 t, 셋 이상 집단의 평균이면 F,
분포 모양 전체면 $D$다.

% =====================================================================
\section{무엇을 비교할지 고르기}
\label{sec:cd-choose}
% source: deck 11 slide 2 (덱 10 슬라이드 22, 23과 같음)

질문 셋을 먼저 본다. 치료가 병 발생률을 낮추는가(평균 또는 비율의
비교), 보스턴 마라톤 여성 주자가 남성보다 젊은가(평균의 비교, 또는
나이 분포 전체의 비교), 어떤 유전자가 전이암에서 다르게 발현되는가(수천
번 반복되는 평균 비교). 셋 모두 "무엇을 비교하는가"라는 같은 질문의
다른 얼굴이다.

그러나 통계량을 고르기 전에 먼저 물어야 할 질문이 하나 더 있다(덱
10 슬라이드 23). 서로 다른 변수를, 서로 다른 모집단에서 재고 있지는
않은가. 비교 자체가 성립하는지부터 물어야 한다. 변수나 모집단이
다르면 어떤 검정을 골라도 그 비교는 의미가 없다.

\Cref{fig:cd-choose}가 그 갈림길을 정리한다.

\begin{figure}[ht]\centering
\begin{tikzpicture}[
  x = 1mm, y = 1mm,
  flowbox/.style = {draw = bookrule, fill = booksurface, rounded corners = 2pt,
                   minimum width = 34mm, minimum height = 9mm, align = center,
                   font = \small, inner sep = 1pt},
  leafbox/.style = {draw = bookaccent, fill = bookaccentbg, rounded corners = 2pt,
                   minimum width = 34mm, minimum height = 11mm, align = center,
                   font = \small, inner sep = 1pt},
  arrow/.style  = {-{Stealth[length=2.2mm]}, draw = bookmuted}]
  \node[flowbox] (q) at (79, 0) {무엇을 비교하는가?};
  \node[flowbox] (mean)  at (40, -18) {평균};
  \node[flowbox] (shape) at (118, -18) {분포 전체};
  \node[leafbox] (t2)  at (20, -36)  {집단 2개\\[-1pt]\footnotesize t검정};
  \node[leafbox] (f)   at (60, -36)  {집단 3개 이상\\[-1pt]\footnotesize ANOVA(F검정)};
  \node[leafbox] (ks1) at (98, -36)  {가설 분포 있음\\[-1pt]\footnotesize 한 표본 KS};
  \node[leafbox] (ks2) at (138, -36) {가설 분포 없음\\[-1pt]\footnotesize 두 표본 KS};
  \draw[arrow] (q) -- (mean);
  \draw[arrow] (q) -- (shape);
  \draw[arrow] (mean) -- (t2);
  \draw[arrow] (mean) -- (f);
  \draw[arrow] (shape) -- (ks1);
  \draw[arrow] (shape) -- (ks2);
\end{tikzpicture}
\caption{무엇을 비교할지 고르는 갈림길. 평균을 비교하면 집단 수에 따라
t검정 또는 ANOVA, 분포 모양 전체를 비교하면 가설 분포의 유무에 따라
한 표본 또는 두 표본 KS검정.}
\label{fig:cd-choose}
\end{figure}

% =====================================================================
\section{단일 표본 t검정}
\label{sec:cd-t-one}
% source: deck 11 slide 3, 4, 5

\begin{inbrief}
  \item[Why] 모집단 표준편차를 모르므로 $z$검정을 그대로 쓸 수 없다.
             표준편차를 표본에서 추정하면 그 추정 오차만큼 분포가
             퍼진다.
  \item[When] 평균 하나를 기준값 $\mu_0$와 비교할 때. 또는 한 집단을
              두 조건에서 잰 차이(\term{대응 설계}(Within-subjects
              Design))를 볼 때.
  \item[Where] 배터리 평균 수명 10시간, 시험 평균 75점, 혈압 감소
               8\,mmHg, 교사 경력, 반바지 대 긴바지 달리기 기록.
  \item[How] $t = (\bar x - \mu_0)/(s/\sqrt n)$을 계산하고, 자유도
             $n-1$인 t분포(t-distribution)에서 p값을 읽는다.
\end{inbrief}

\paragraph{고치는 문제.} $z$통계량(\Cref{sec:pv-continuous})은
$z = (\bar x - \mu_0)/(\sigma/\sqrt n)$처럼 모집단 표준편차
$\sigma$를 쓴다. 그런데 실제 자료에서는 $\sigma$를 모른다. 표본
표준편차 $s$로 대신 쓰면, $s$ 자체가 표본마다 흔들리는 추정값이라서
그 흔들림만큼 분포가 정규분포보다 더 퍼진다.

\paragraph{아이디어.} 분모의 $\sigma$를 $s$로 바꾸고, 그 대체가
만드는 추가 흔들림을 반영하도록 정규분포 대신 꼬리가 더 두꺼운
t분포를 쓴다. 표본이 작을수록($n$이 작을수록) $s$의 추정 오차가
크므로 꼬리가 더 두껍고, 표본이 커지면 $s \to \sigma$에 가까워지며
t분포는 정규분포에 수렴한다.

\begin{definition}[t통계량]
\label{def:t-statistic}
단일 표본에서 평균 $\bar x$가 기준값 $\mu_0$와 다른지 검정할 때는

\[
  t = \frac{\bar x - \mu_0}{s/\sqrt n}
\]

를 쓴다. 자유도는 $n-1$이다.

두 독립 표본의 평균 $\bar x_1$, $\bar x_2$가 다른지 검정할 때는

\[
  t = \frac{\bar x_1 - \bar x_2}{\sqrt{s_1^2/n_1 + s_2^2/n_2}}
\]

를 쓴다(\Cref{sec:cd-t-two}). 자유도는 계산 방식에 따라 달라진다
(\Cref{sec:cd-t-two}의 함정 참고).
\end{definition}

\begin{pitfall}[슬라이드 값은 $t = \bar x/\bar\sigma$, 다시 계산하면
$t = (\bar x - \mu_0)/(s/\sqrt n)$]
슬라이드 3(\texttt{image1})의 식은 $\mu_0$와 $\sqrt n$이 모두 빠져
있다. 같은 슬라이드의 \texttt{image11}과 슬라이드 5의 \texttt{image5}는
올바른 식을 보여준다. $\mu_0$가 없으면 "기준값과 다른가"를 물을 수
없다. 대응 설계처럼 $\mu_0 = 0$인 경우는 분자만 보면 우연히 맞아
보이지만, 문제는 분모에 있다. $\sqrt n$을 나누는 이유는 표본이
커질수록 평균 $\bar x$의 표준오차 $s/\sqrt n$이 작아지고, 같은 크기의
차이 $\bar x - \mu_0$라도 표준오차 단위로는 더 크게(더 확실하게)
잡히기 때문이다. $t$는 그렇게 표본이 늘수록 커져서 증거를 더 강하게
반영해야 한다. $\sqrt n$을 빼면 분모가 $s$(표본 표준편차) 하나로
고정되어 $n$과 무관해지고, 그러면 $t$는 표본이 아무리 커져도
커지지 않는다. 즉 데이터를 더 모아도 통계량에 그 증거가 반영되지
않는다.
\end{pitfall}

슬라이드 4는 계산 없이 예시 셋을 든다. 배터리 공장은 평균 수명이
10시간인지, 학교는 평균 점수가 75점인지, 제약회사는 평균 혈압 감소가
8\,mmHg인지 각각 단일 표본 t검정으로 확인한다.

\begin{example}[교사 경력]
\label{ex:cd-rory}
슬라이드 5의 값이다. $n = 25$, $\bar x = 4$, $s = 2$이고, 대립가설은
$H_a: \mu < 5$(단측)다.

\begin{align*}
  &s/\sqrt n = 2/\sqrt{25} = 2/5 = 0.4 && \text{표준오차}\\
  &t = (4 - 5)/0.4 = -2.5 && \text{t통계량}\\
  &\text{df} = n - 1 = 24 && \text{자유도}\\
  &P(t_{24} < -2.5) \approx 0.0098 && \text{단측 p값(\texttt{t.cdf(-2.5, 24)})}
\end{align*}

$\alpha = 0.05$ 기준으로 $0.0098 < 0.05$이므로 기각한다. 슬라이드는
결론을 밝히지 않으므로, 이 결론 문장은 \Cref{tab:ht-wording}의 형식을
따른 것이다.
\end{example}

\Cref{fig:cd-t-dist}는 $N(0,1)$과 자유도 5, 24인 t분포를 겹쳐 그려서
"자유도가 작을수록 꼬리가 두껍다, 같은 $t$값이라도 p값이 더 크게
나온다"를 보여준다.

\begin{figure}[ht]\centering
  \includegraphics[width=0.85\linewidth]{11_01_t_vs_normal.png}
  \caption{$N(0,1)$과 t분포(자유도 5, 24) 비교. 자유도 5 곡선의
  $\pm 2.75$ 바깥 꼬리를 칠해 \Cref{ex:cd-t-table}의 양측 p값
  $0.04$를 표시했다. 꼬리가 두꺼울수록 같은 $t$에서 p값이 크다.}
  \label{fig:cd-t-dist}
\end{figure}

% =====================================================================
\section{t표 읽기}
\label{sec:cd-t-table}
% source: deck 11 slide 9

\begin{inbrief}
  \item[Why] 소프트웨어가 없을 때도 p값의 범위를 알아야 한다.
  \item[When] $t$값과 자유도를 구한 뒤.
  \item[Where] $t = 2.75$, $n = 6$.
  \item[How] 자유도 행에서 $t$값을 끼는 두 열을 찾고, 그 열의 꼬리
             넓이를 읽는다. 양측검정이면 두 배로 만든다.
\end{inbrief}

\paragraph{고치는 문제.} \Cref{def:t-statistic}로 $t$값을 구해도,
소프트웨어 없이 손으로는 그 $t$가 만드는 정확한 꼬리 넓이(p값)를
적분으로 구할 수 없다. 강의 당시(그리고 시험장처럼 소프트웨어가
없는 상황)에는 미리 계산해 둔 표를 봐야 한다.

\paragraph{아이디어.} t표는 자유도별로 몇 개의 대표적인 꼬리 넓이
(예: $0.10$, $0.05$, $0.02$, $0.01$)에 해당하는 임계값만 미리
계산해 둔다. 구한 $t$값이 표의 어느 두 임계값 사이에 있는지 찾으면,
p값이 그 두 꼬리 넓이 사이에 있다는 것을 알 수 있다. 정확한 값은
아니지만 기각 여부를 판단하기에는 충분하다.

\begin{example}[t표 읽기]
\label{ex:cd-t-table}
$t = 2.75$, $n = 6$이므로 자유도는 $n - 1 = 5$다. 자유도 5 행에서
$2.75$는 한쪽 꼬리 넓이 $0.02$에 해당하는 임계값 $t_{0.02,5} = 2.757$
바로 아래에 있다.

\begin{align*}
  &P(t_5 > 2.75) \approx 0.02 && \text{한쪽 꼬리(표에서 읽은 값)}\\
  &2 \times 0.02 = 0.04 && \text{양측(부호를 모르므로 두 배)}\\
  &2 \times P(t_5 > 2.75) = 0.0403 && \text{정확 계산(\texttt{2*t.sf(2.75,5)})}
\end{align*}

슬라이드의 $\text{p-value} = 0.04$는 정확 계산 $0.0403$과 반올림
범위 안에서 일치한다. 슬라이드가 t표 전체(\texttt{image20})와 일부
행(\texttt{image21}, \texttt{image22})을 보여주지만, 이 책은 표를
그대로 옮기지 않고 필요한 셀 $2.757$ 하나만 인용한다.
\end{example}

% =====================================================================
\section{두 표본 t검정}
\label{sec:cd-t-two}
% source: deck 11 slide 6, 7, 8

\begin{inbrief}
  \item[Why] 서로 다른 사람들로 이뤄진 두 집단(\term{독립 설계}(Between-
             subjects Design))은 대응 설계처럼 차이를 사람별로
             뺄 수 없다.
  \item[When] 두 독립 표본의 평균을 비교할 때.
  \item[Where] 남녀 급여, 과외 유무에 따른 성적, 두 약의 혈압 감소,
               두 매장의 매출, 두 밭의 토마토 키.
  \item[How] $t = (\bar x_1 - \bar x_2)/\sqrt{s_1^2/n_1 + s_2^2/n_2}$을
             계산하고, 자유도를 정해 p값을 읽는다.
\end{inbrief}

\paragraph{고치는 문제.} \Cref{sec:cd-t-one}의 단일 표본 $t$는 표본
하나를 고정된 기준값 $\mu_0$와 비교한다. 그런데 남녀 급여, 두 밭의
토마토처럼 비교할 대상이 \textbf{둘 다 표본}이면, 기준값 자체가 없고
두 표본 평균의 차이 $\bar x_1 - \bar x_2$를 재야 한다. 이 차이에는
두 표본 모두의 흔들림(각각의 표준오차)이 섞여 들어간다.

\paragraph{아이디어.} 분모의 표준오차를 표본 하나짜리 $s/\sqrt n$에서
두 표본의 분산을 각각 표본 크기로 나눠 더한 값의 제곱근
$\sqrt{s_1^2/n_1 + s_2^2/n_2}$로 바꾼다. 두 표본이 독립이면 분산은
더해진다는 성질, 즉 $\mathrm{Var}(\bar x_1 - \bar x_2) = \sigma_1^2/n_1 + \sigma_2^2/n_2$
를 그대로 쓴 것이다. 빼도 분산은 줄지 않고 더해진다. 두 평균이 각자 흔들리고, 그 흔들림이 서로 상쇄될 이유가 없기 때문이다. 이렇게 만든
두 표본 t통계량이 \Cref{def:t-statistic}의 두 번째 식이다.

슬라이드 7은 계산 없이 예시 넷을 든다. 과외 유무에 따른 시험 점수,
두 약의 혈압 감소, 두 매장의 매출, 남녀 운동선수의 심박수가 그것이다.

\begin{example}[두 밭의 토마토]
\label{ex:cd-tomato}
슬라이드 8(덱 10 슬라이드 25와 같음)의 값이다. \Cref{tab:cd-tomato}가
두 밭 A, B의 평균, 표준편차, 표본 크기를 정리한다.

\begin{table}[ht]\centering
\begin{tabularx}{\linewidth}{@{}l X X X X@{}}
\toprule
밭 & 평균(Mean) $\bar x$ & 표준편차(SD) $s$ & $n$ & $s^2/n$ \\
\midrule
A & 1.3 & 0.5 & 22 & 0.011364 \\
B & 1.6 & 0.3 & 24 & 0.00375 \\
\bottomrule
\end{tabularx}
\caption{두 밭의 토마토 키 요약(단위 생략).}
\label{tab:cd-tomato}
\end{table}

\begin{align*}
  &s_A^2/n_A = 0.5^2/22 = 0.011364 && \text{A의 분산/표본크기}\\
  &s_B^2/n_B = 0.3^2/24 = 0.00375 && \text{B의 분산/표본크기}\\
  &0.011364 + 0.00375 = 0.015114 && \text{합}\\
  &\sqrt{0.015114} = 0.12294 && \text{표준오차}\\
  &t = (1.3 - 1.6)/0.12294 = -2.4403 && \text{t통계량}
\end{align*}

$|t| \approx 2.44$다. 슬라이드는 $P(|t| \ge 2.44) \approx 0.024$라고
적는다.
\end{example}

\begin{pitfall}[슬라이드 값은 $p \approx 0.024$(자유도 미표기), 다시
계산하면 자유도를 무엇으로 잡느냐에 따라 $p$가 달라진다]
슬라이드는 자유도를 밝히지 않는다. 보수적 방식(작은 쪽 표본 크기
$-1$, 즉 $\min(n_A, n_B) - 1 = 21$)을 쓰면

\[
  2 \times P(t_{21} > 2.44) = 0.0236 \approx 0.024
\]

로 슬라이드 값과 일치한다. 그러나 소프트웨어 기본값(Welch의 근사
자유도, \texttt{scipy.stats.ttest\_ind\_from\_stats(..., equal\_var=False)})은

\[
  \text{df} \approx 33.79 \approx 33.8, \qquad p \approx 0.0201 \approx 0.020
\]

을 준다. 자유도가 커질수록 p값이 조금 작아지지만, 결론은 둘 다 같다.
$\alpha = 0.05$에서 두 값 모두 기각이다. 결론이 $\alpha$ 경계 근처에
있는 계산이라면 어느 자유도를 썼는지 밝혀야 한다. 이 책의 본문 공식은
슬라이드 그대로 쓰고, p값은 두 방식을 나란히 적는다.
\end{pitfall}

% =====================================================================
\section{분산분석과 F검정}
\label{sec:cd-anova}
% source: deck 11 slide 10, 11, 12

\begin{inbrief}
  \item[Why] 집단이 셋 이상이면 t검정을 쌍마다 반복해야 한다. 검정을
             여러 번 하면 거짓 양성이 쌓인다.
  \item[When] 셋 이상 집단의 \textbf{평균}이 모두 같은지 한 번에
              물을 때.
  \item[Where] 대학생, 고등학생, 초등학생의 만화 시청 시간.
  \item[How] $F = \dfrac{\mathrm{SSB}/(m-1)}{\mathrm{SSW}/(m(n-1))}$을
             계산한다. 값이 크면 기각한다. F표에서 임계값을 읽는다.
\end{inbrief}

\paragraph{고치는 문제.} 집단이 셋 이상이면 두 표본 t검정을 쌍마다
반복하면 된다고 생각하기 쉽다. 그러나 검정을 여러 번 반복하면
\Cref{sec:pv-threshold}의 문제, 즉 매 검정마다 $5\%$씩 쌓이는 거짓
양성이 그대로 발생한다. 집단이 $m$개면 쌍의 수는
$\binom{m}{2}$개이므로, $m$이 커질수록 거짓 양성이 빠르게 늘어난다.

\paragraph{아이디어.} 쌍마다 비교하는 대신, "집단 사이의 흩어짐"과
"집단 안의 흩어짐"을 한 번에 비교하는 통계량을 만든다. 집단
평균들이 서로 멀리 떨어져 있으면(흩어짐이 크면) 평균이 다르다는
증거이고, 그 흩어짐을 같은 집단 안의 자연스러운 흩어짐과 비율로
비교하면 우연인지 아닌지 판단할 수 있다. 이 비율이 $F$통계량이다.

\begin{definition}[F통계량]
\label{def:f-statistic}
집단이 $m$개, 집단마다 표본 크기가 $n$일 때

\[
  F = \frac{\mathrm{SSB}/(m-1)}{\mathrm{SSW}/(m(n-1))}
\]

를 쓴다. 분자는 \term{집단 간 제곱합}(SSB)을 그 자유도 $m-1$로 나눈
값, 분모는 \term{집단 내 제곱합}(SSW)을 그 자유도 $m(n-1)$로 나눈
값이다. $F$가 크다는 것은 집단 사이의 흩어짐이 집단 안의 흩어짐보다
훨씬 크다는 뜻이고, 이는 집단 평균이 서로 다르다는 증거다.
\end{definition}

\begin{pitfall}[슬라이드는 "Comparing Variances", 바르게는 "평균을
비교하는 도구가 분산 비율이다"]
슬라이드 2는 "Variances $\to$ F-test"라 적고, 슬라이드 10의 제목은
"Comparing Variances"다. 그러나 ANOVA의 귀무가설은 "모든 집단의
\textbf{평균}이 같다"이다. 분산 비율($F$)은 그 판정에 쓰는
\textbf{통계량}일 뿐, ANOVA가 비교하는 대상은 평균이다. 두 분산 자체가
같은지 묻는 검정도 같은 F분포(F-distribution)를 쓰지만, 그것은
\textbf{다른} 검정이다.
이 덱에는 그 검정의 예제가 없으므로 이름만 적고 계산하지 않는다.
\end{pitfall}

\begin{keyidea}[핵심: 집단이 둘이면 $F = t^2$]
집단이 둘뿐이면 ANOVA와 두 표본 t검정(\term{합동 분산}(Pooled
Variance) 가정)은 같은 질문을
다른 통계량으로 묻는 것이고, $F = t^2$가 성립한다. \Cref{ex:cd-revisit}가
이것을 확인한다.
\end{keyidea}

\begin{example}[만화 시청 시간]
\label{ex:cd-cartoon}
대학생, 고등학생, 초등학생 각 3명이 만화를 본 시간이다. 대학생
$\{3, 2, 1\}$, 고등학생 $\{5, 3, 4\}$, 초등학생 $\{5, 6, 7\}$.

\begin{align*}
  &(3+2+1)+(5+3+4)+(5+6+7) = 36 && \text{전체 합}\\
  &36/9 = 4 && \text{전체 평균}\\
  &2,\ 4,\ 6 && \text{집단 평균(대학생, 고등학생, 초등학생)}
\end{align*}

\begin{align*}
  &(3{-}2)^2{+}(2{-}2)^2{+}(1{-}2)^2 = 2 && \text{대학생 내부 편차 제곱합}\\
  &(5{-}4)^2{+}(3{-}4)^2{+}(4{-}4)^2 = 2 && \text{고등학생 내부}\\
  &(5{-}6)^2{+}(6{-}6)^2{+}(7{-}6)^2 = 2 && \text{초등학생 내부}\\
  &2+2+2 = 6 && \text{SSW(집단 내 제곱합)}\\
  &(2{-}4)^2+(4{-}4)^2+(6{-}4)^2 = 8 && \text{집단 평균의 편차 제곱합}\\
  &3 \times 8 = 24 && \text{SSB(집단마다 3명이므로 3배)}\\
  &24 + 6 = 30 && \text{SST(전체 제곱합, 검산)}
\end{align*}

검산: $\sum (x-4)^2 = 1+4+9+1+1+0+1+4+9 = 30$으로 위 합과 일치한다.

\begin{align*}
  &\text{df}_{\text{SSB}} = m - 1 = 2 && \text{집단 수} - 1\\
  &\text{df}_{\text{SSW}} = m(n-1) = 3 \times 2 = 6 && \text{집단마다 자유도 2, 3집단}\\
  &\text{df}_{\text{SST}} = mn - 1 = 8 = 2 + 6 && \text{검산}\\
  &F = \frac{24/2}{6/6} && \text{F통계량의 정의대로 대입}\\
  &\phantom{F} = \frac{12}{1} = 12 && \text{분자와 분모를 계산}
\end{align*}

\Cref{tab:cd-anova}가 이 계산을 분산분석표로 정리한다.

\begin{table}[ht]\centering
\begin{tabularx}{\linewidth}{@{}l X X X X@{}}
\toprule
요인 & 제곱합(SS) & 자유도(df) & 평균제곱(Mean Square, MS) & $F$ \\
\midrule
집단 간(Between) & 24 & 2 & 12 & 12 \\
집단 내(Within) & 6 & 6 & 1 & -- \\
\midrule
전체(Total) & 30 & 8 & -- & -- \\
\bottomrule
\end{tabularx}
\caption{만화 시청 시간 분산분석표(ANOVA Table).}
\label{tab:cd-anova}
\end{table}

$\alpha = 0.10$에서 임계값 $F_{0.10}(2,6) = 3.46$이고, $12 > 3.46$이므로
기각한다. 학년(교육 수준)에 따라 만화 시청 시간의 평균이 다르다.
\Cref{fig:cd-anova}가 SSW와 SSB 두 거리를 점그림으로 보여준다.

\begin{figure}[ht]\centering
  \includegraphics[width=0.85\linewidth]{11_02_anova_cartoon.png}
  \caption{만화 시청 시간 예제의 점그림. 왼쪽: 각 점에서 자기 집단
  평균까지의 거리(SSW). 오른쪽: 각 집단 평균에서 전체 평균 4까지의
  거리(SSB). 집단 평균선 2, 4, 6과 전체 평균선 4를 함께 표시했다.}
  \label{fig:cd-anova}
\end{figure}
\end{example}

\begin{pitfall}[슬라이드 값은 SSW의 자유도가 $m*n-1$, 다시 계산하면
$m(n-1)$]
슬라이드 11은 대학생($2$), 고등학생($2$), 초등학생($2$) 세 집단의
내부 자유도를 $2+2+2=6$으로 더한 뒤, 그 결과 $6$ 옆에
"$m*n-1$"이라는 라벨을 붙인다. 그러나 $m \times n - 1 = 3 \times 3 -
1 = 8$은 \textbf{전체} 자유도이고, SSW의 자유도는 각 집단의 자유도
$n-1$을 $m$개 더한 $m(n-1) = 3 \times 2 = 6$이다. 라벨의 식
$m*n-1=8$과 실제로 계산에 쓴 값 $6$이 다른데도, 슬라이드는 같은
숫자 $6$ 옆에 틀린 라벨을 붙여 둘이 같은 것처럼 보이게 만든다.
\Cref{tab:cd-anova}처럼 세 줄로 나누면 이 혼동이 사라진다.
\end{pitfall}

\begin{pitfall}[슬라이드 값은 "$\alpha < 0.10$", 다시 계산하면
$\alpha = 0.10$]
슬라이드 11(\texttt{image31})은 유의수준을 "significance level
$\alpha < 0.10$"이라 적는다. 유의수준은 부등식이 아니라 \textbf{정한
값}이다. 슬라이드 12는 실제로 $\alpha = 0.10$으로 맞게 쓴다.
\end{pitfall}

임계값 $F_{0.10}(2,6) = 3.46$은 F표에서 분자 자유도 2, 분모 자유도 6
칸을 읽은 값이다(\texttt{f.ppf(0.90, 2, 6)} $= 3.4633$). 분자
자유도가 2인 F분포는 꼬리 확률을 닫힌 식으로 쓸 수 있다.

\begin{align*}
  &P(F > f) = \left(1 + \frac{2f}{\nu_2}\right)^{-\nu_2/2}
    && \text{분자 자유도 2 전용 공식}\\
  &\left(1 + \frac{24}{6}\right)^{-3} = 5^{-3} = 0.008
    && \text{$f=12$, $\nu_2=6$ 대입}
\end{align*}

$p = 0.008$이므로 $\alpha = 0.01$ 기준으로도 기각이다.

\begin{example}[약 C 대 D 다시 보기]
\label{ex:cd-revisit}
\Cref{ex:ht-distance}는 $C = \{20, 28, 34\}$, $D = \{22, 26,
37\}$에서 합친 평균 대 집단 평균의 제곱거리 합만 비교하고, "정식
검정 통계량으로 만드는 방법은 이 덱에 없다"며 멈췄다. 이제 ANOVA와
두 표본 t검정으로 그 질문을 닫는다(이 값은 슬라이드가 아니라 이
책 자체의 예시다).

\Cref{ex:ht-distance}가 이미 계산한 값을 그대로 쓴다. 집단 간
제곱거리 합(줄어든 양) $\text{SSB} = 1.50$, 집단 내 제곱거리 합
$\text{SSW} = 219.33$이다. 집단이 둘($m=2$), 집단마다 $n=3$이므로
자유도는 $\text{df}_{\text{SSB}} = m - 1 = 1$,
$\text{df}_{\text{SSW}} = m(n-1) = 4$다.

\begin{align*}
  &F = \frac{1.50/1}{219.33/4} && \text{F통계량(\Cref{def:f-statistic})에 대입}\\
  &\phantom{F} = \frac{1.50}{54.83} && \text{분모를 계산}\\
  &\phantom{F} = 0.0274 && \text{나눗셈}
\end{align*}

같은 데이터를 두 표본 t검정으로도 계산한다. $\bar x_D - \bar x_C =
28.33 - 27.33 = 1$이고, 합동 분산은 SSW를 합동 자유도로 나눈 값
$219.33/4 = 54.83$이다.

\begin{align*}
  &54.83 \times (1/3 + 1/3) = 36.55 && \text{합동 분산} \times \text{표준오차 계수}\\
  &\sqrt{36.55} = 6.046 && \text{표준오차}\\
  &t = 1/6.046 = 0.165 && \text{t통계량}\\
  &t^2 = 0.0274 = F && \text{검산(\Cref{def:t-statistic}, \Cref{def:f-statistic})}\\
  &2 \times P(t_4 > 0.165) \approx 0.88 && \text{양측 p값}
\end{align*}

$p \approx 0.88$은 $\alpha = 0.05$는커녕 어떤 통상적인 기준으로도
기각할 수 없다. 덱 9의 "극적인 차이가 없다"는 결론이 이제 수치로
확인된다. 이 예제가 \Cref{sec:ht-future}가 열어 둔 질문의 답이다.
\end{example}

% =====================================================================
\section{KS 검정: 표본 대 가설 분포}
\label{sec:cd-ks}
% source: deck 11 slide 13, 14, 15, 16, 17, 18, 19, 20

\begin{inbrief}
  \item[Why] 평균과 분산이 같아도 분포의 모양 전체는 다를 수 있다.
  \item[When] 표본이 특정 분포(예: $N(0,1)$)에서 왔는지 물을 때.
  \item[Where] 관측값 100개가 표준정규분포에서 왔는가.
  \item[How] 데이터를 정렬한다 $\to$ 경험적 누적분포함수
             $F_{\text{obs}}$를 만든다 $\to$ 가설 분포의 누적분포함수
             $F_{\text{exp}}$와 비교한다 $\to$ 세로 격차의 최댓값
             $D_n$을 구한다.
\end{inbrief}

\paragraph{고치는 문제.} $t$검정과 F검정은 \textbf{평균}만 비교한다.
그런데 평균과 분산이 완전히 같아도, 한쪽은 좌우 대칭이고 다른 쪽은
한쪽으로 치우친 것처럼 분포의 \textbf{모양 전체}가 다를 수 있다.
평균과 분산만 보는 검정으로는 이런 차이를 잡아내지 못한다.

\paragraph{아이디어.} 모양 전체를 비교하려면 분포 전체를 나타내는
함수, 즉 누적분포함수(\Cref{sec:cdf})를 직접 비교하면 된다. 두
누적분포함수를 겹쳐 그리고, 가장 크게 벌어지는 한 지점의 세로 거리를
잰다. 이것이 콜모고로프-스미르노프 검정(KS 검정, Kolmogorov-Smirnov
Test)이다(슬라이드 13). "두 지도를 겹쳤을 때 가장 어긋나는 한 지점의
거리"라는 비유가 그대로 정의다. \Cref{fig:cd-ks-one}이 그 거리를
그림으로 보여준다(슬라이드 14, 20이 같은 그림을 반복한다. 이 책은
한 번만 그린다).

\begin{figure}[ht]\centering
  \includegraphics[width=0.85\linewidth]{11_03_ks_one_sample.png}
  \caption{$N(0,1)$에서 뽑았다고 주장하는 표본 100개의 경험적
  누적분포함수(Empirical CDF, 계단)와 이론적 누적분포함수(Theoretical
  CDF) $\Phi$(매끄러운 곡선).
  27번째 값에서 계단 \textbf{아래}로 잰 격차 $0.096$과, 88번째
  값에서 계단 \textbf{위}로 잰 격차 $0.092$를 함께 표시했다.}
  \label{fig:cd-ks-one}
\end{figure}

\begin{definition}[경험적 누적분포함수]
\label{def:ecdf}
표본 $x_1, \dots, x_n$을 오름차순으로 정렬했을 때, 값 $x$에서의
경험적 누적분포함수는

\[
  F_{\text{obs}}(x) = \frac{\#\{i : x_i \le x\}}{n}
\]

다. 정렬된 $i$번째 값에서는 $F_{\text{obs}}(x_i) = i/n$이 된다.
\end{definition}

\begin{pitfall}[슬라이드 값은 "$x$ 미만(below)의 개수", 바르게는
"$x$ 이하(at or below)의 개수"]
슬라이드 15(\texttt{image34})는 $F_{\text{obs}}(x)$를 "\#observations
below $x$"로 적는다. 그러나 $i$번째 정렬값에서 $F_{\text{obs}}(x_i) =
i/n$이 성립하려면, $x_i$ 자신도 세어야 한다. 그래서 \Cref{def:ecdf}는
"이하"로 정의한다.
\end{pitfall}

가설 분포의 누적분포함수 $F_{\text{exp}}(x_i) = \Phi(x_i)$는 표준정규
누적분포함수 표에서 읽는다(슬라이드 18).

\begin{definition}[KS 통계량]
\label{def:ks-statistic}
정렬된 표본에서 관측 누적분포함수 $F_{\text{obs}}$와 기대 누적분포함수
$F_{\text{exp}}$ 사이의 최대 세로 격차를

\[
  D_n = \sup_x \left| F_{\text{obs}}(x) - F_{\text{exp}}(x) \right|
\]

로 정의한다. 표본이 실제로 그 가설 분포에서 왔다면 $D_n$은 작아야
한다.
\end{definition}

\begin{aside}[참고]
표본이 실제로 가설 분포에서 왔다면, 표본 크기 $n \to \infty$일 때
$D_n \to 0$으로 수렴한다는 사실이 증명되어 있다(슬라이드 16). 이
사실을 \term{글리벤코-칸텔리 정리}(Glivenko-Cantelli Theorem)라
부른다. 슬라이드는 이 이름을 쓰지 않는다.
\end{aside}

\begin{example}[100개가 $N(0,1)$인가]
\label{ex:cd-ks-normal}
슬라이드 17~19의 예제다. 관측값 100개(\texttt{image40})를 정렬하고,
정렬 결과가 \texttt{image39}와 같은지 확인한 뒤(전사 검증), $i$번째
값에서 $F_{\text{obs}}(x_i) = i/100$과 $\Phi(x_i)$를 비교한다. 표
전체를 옮기지 않고 두 행만 본다.

\begin{align*}
  &i = 27,\ x_{27} = -0.37 && \text{27번째 값}\\
  &\Phi(-0.37) = 0.356 && \text{기대 누적분포함수}\\
  &(i-1)/100 = 26/100 = 0.260 && \text{계단 바로 아래}\\
  &|0.356 - 0.260| = 0.096 && \text{아래쪽 격차}
\end{align*}

\begin{align*}
  &i = 88,\ x_{88} = 0.80 && \text{88번째 값}\\
  &\Phi(0.80) = 0.7881 && \text{기대 누적분포함수}\\
  &i/100 = 88/100 = 0.880 && \text{계단 위}\\
  &|0.880 - 0.7881| = 0.092 && \text{위쪽 격차}
\end{align*}

두 값 중 큰 쪽이 $D_n$이다. \texttt{scipy.stats.kstest(data, 'norm')}로
확인하면 $D_n = 0.0957 \approx 0.096$, $p \approx 0.2997$다.
\end{example}

\begin{pitfall}[슬라이드 값은 $D_n = 0.092$, 다시 계산하면 $D_n =
0.096$]
슬라이드 19는 $i$번째 값에서 $i/100 - \Phi(x_i)$, 즉 계단
\textbf{위쪽} 격차만 계산해 $0.092$를 얻는다. 그러나
\Cref{def:ks-statistic}의 $D_n$은 계단 \textbf{양쪽}(위와 아래)에서
잰 격차 중 큰 값이다. 27번째 값의 아래쪽 격차 $0.096$이 88번째 값의
위쪽 격차 $0.092$보다 크므로, 정확한 $D_n \approx 0.096$이다.
결론은 같다. 임계값 $1.36/\sqrt{100} = 0.136$보다 작으므로 기각
실패다.
\end{pitfall}

임계값 $0.136$은 표본이 클 때 쓰는 근사식 $1.36/\sqrt n$($\alpha =
0.05$, 한 표본)에서 나온다. 슬라이드 23의 임계값표(\texttt{image45})는
$n \le 6$까지만 보이므로 $n = 100$인 이 예제에는 직접 쓸 수 없다.
큰 표본에는 근사식을 쓴다(\Cref{sec:cd-ks-pvalue}에서 이어진다).

모수를 표본 자체에서 추정한 뒤 KS검정을 쓰면 p값이 낙관적으로
나온다(\Cref{sec:normal}의 함정). 이 예제는 $N(0,1)$을
\textbf{미리} 정했으므로 해당하지 않지만, HW2처럼 모수를 데이터에서
추정한 경우는 이 함정에 해당한다.

% =====================================================================
\section{KS 검정: 두 표본}
\label{sec:cd-ks-two}
% source: deck 11 slide 21, 22

\begin{inbrief}
  \item[Why] 가설 분포가 없고 표본만 둘 있을 때도 모양을 비교해야
             한다.
  \item[When] 두 표본이 같은 분포에서 왔는지 물을 때.
  \item[Where] $X$ 10개 대 $Y$ 8개, HW2의 데이터 대 모형 표본.
  \item[How] 합친 값을 정렬하며 두 경험적 누적분포함수를 같이
             올리고, 가장 큰 격차를 $D$로 잡는다.
\end{inbrief}

두 표본에서는 가설 분포가 없으므로 $F_{\text{exp}}$ 대신 두 번째
표본의 경험적 누적분포함수를 쓴다. \Cref{def:ks-statistic}의 $D_n$
정의가 그대로 적용되고, 다만 $F_{\text{exp}}$ 자리에 $Y$의 경험적
누적분포함수가 들어간다.

\begin{example}[표본 $X$ 대 $Y$]
\label{ex:cd-ks-two}
슬라이드 21(제목 "Comparing to Samples"는 "two"의 오타로 보인다)의
값이다. $X$는 10개, $Y$는 8개다.

\[
  X = \{1.2, 1.4, 1.9, 3.7, 4.4, 4.8, 9.7, 17.3, 21.1, 28.4\}
\]
\[
  Y = \{5.6, 6.5, 6.6, 6.9, 9.2, 10.4, 10.6, 19.3\}
\]

$X$는 계단 높이 $1/10$, $Y$는 계단 높이 $1/8$이다.
\Cref{tab:cd-ks-two}가 합친 18개 값을 모두 정렬해 각 지점의
$F_x$, $F_y$, $|F_x - F_y|$를 보여준다.

\begin{table}[ht]\centering
\begin{tabularx}{\linewidth}{@{}l X X X@{}}
\toprule
값 & $F_x$ & $F_y$ & $|F_x - F_y|$ \\
\midrule
1.2 & 0.1 & 0 & 0.1 \\
1.4 & 0.2 & 0 & 0.2 \\
1.9 & 0.3 & 0 & 0.3 \\
3.7 & 0.4 & 0 & 0.4 \\
4.4 & 0.5 & 0 & 0.5 \\
4.8 & 0.6 & 0 & \textbf{0.6} \\
5.6 & 0.6 & 0.125 & 0.475 \\
6.5 & 0.6 & 0.25 & 0.35 \\
6.6 & 0.6 & 0.375 & 0.225 \\
6.9 & 0.6 & 0.5 & 0.1 \\
9.2 & 0.6 & 0.625 & 0.025 \\
9.7 & 0.7 & 0.625 & 0.075 \\
10.4 & 0.7 & 0.75 & 0.05 \\
10.6 & 0.7 & 0.875 & 0.175 \\
17.3 & 0.8 & 0.875 & 0.075 \\
19.3 & 0.8 & 1.0 & 0.2 \\
21.1 & 0.9 & 1.0 & 0.1 \\
28.4 & 1.0 & 1.0 & 0 \\
\bottomrule
\end{tabularx}
\caption{$X$, $Y$를 합쳐 정렬한 표 전체(18개 값). 최댓값 행을
강조했다.}
\label{tab:cd-ks-two}
\end{table}

\begin{align*}
  &F_x(4.8) = 6/10 = 0.6 && \text{$X$의 6개(1.2..4.8)가 4.8 이하}\\
  &F_y(4.8) = 0/8 = 0 && \text{$Y$ 원소 중 4.8 이하는 없음}\\
  &|0.6 - 0| = 0.6 && \text{4.8에서의 격차}
\end{align*}

\begin{align*}
  &F_x(5.6) = 6/10 = 0.6 && \text{$X$의 6개가 5.6 이하}\\
  &F_y(5.6) = 1/8 = 0.125 && \text{$Y$의 1개(5.6)가 5.6 이하}\\
  &|0.6 - 0.125| = 0.475 && \text{5.6에서의 격차(다음 후보)}
\end{align*}

\begin{align*}
  &F_x(17.3) = 8/10 = 0.8 && \text{$X$의 8개가 17.3 이하}\\
  &F_y(17.3) = 7/8 = 0.875 && \text{$Y$의 7개가 17.3 이하}\\
  &|0.8 - 0.875| = 0.075 && \text{17.3에서의 격차}
\end{align*}

$0.6 > 0.475 > 0.075$이므로 최댓값은 $D = 0.6$이며, $4.8$에서
일어난다. \Cref{tab:cd-ks-two}의 나머지 행도 모두 $0.6$보다
작으므로 이 결론은 바뀌지 않는다.
\end{example}

\begin{pitfall}[슬라이드 값은 $F_y$가 $0.1$ 단위, 다시 계산하면
$1/8$ 단위]
슬라이드 21의 $F_y$ 열은 $0.1, 0.2, 0.4, 0.5, 0.6, 0.8, 0.9, 1.0$처럼
$0.1$ 단위로 반올림되어 있다. 그러나 $Y$는 8개이므로 계단 한 칸의
높이는 정확히 $1/8 = 0.125$다: $5.6$에서 $0.125$, $6.5$에서
$0.25$, $6.6$에서 $0.375$, $6.9$에서 $0.5$, $9.2$에서 $0.625$,
$10.4$에서 $0.75$, $10.6$에서 $0.875$, $19.3$에서 $1$. 슬라이드
22(\texttt{image44})의 계단 그림은 $1/8$ 단위로 올바르게 그려져
있어, 반올림은 표에만 있는 문제다. $D = 0.6$이라는 결론은 반올림
전후로 바뀌지 않는다.
\end{pitfall}

\Cref{fig:cd-ks-two}가 두 경험적 누적분포함수와 $D = 0.6$이 일어나는
지점을 함께 그린다.

\begin{figure}[ht]\centering
  \includegraphics[width=0.85\linewidth]{11_04_ks_two_sample.png}
  \caption{$X$(계단 높이 $1/10$)와 $Y$(계단 높이 $1/8$)의 경험적
  누적분포함수. $x = 4.8$에서 $F_x = 0.6$, $F_y = 0$으로 격차가
  최대인 $D = 0.6$이다.}
  \label{fig:cd-ks-two}
\end{figure}

% =====================================================================
\section{$D$에서 p값으로}
\label{sec:cd-ks-pvalue}
% source: deck 11 slide 23

\begin{inbrief}
  \item[Why] $D$ 값만으로는 "크다"를 말할 수 없다(슬라이드 23).
  \item[When] $D$를 구한 뒤.
  \item[Where] \Cref{ex:cd-ks-normal}, \Cref{ex:cd-ks-two}.
  \item[How] 임계값표(작은 $n$) 또는 큰 표본 근사식으로 p값을 얻는다.
             실무에서는 \texttt{scipy.stats}의 \texttt{kstest},
             \texttt{ks\_2samp}를 쓴다.
\end{inbrief}

KS 통계량 자체의 분포는 손으로 다루기 어렵지만(슬라이드 23), 큰
표본에서는 근사식이 있다. 한 표본, $\alpha = 0.05$에서는 $1.36/\sqrt
n$을 임계값으로 쓴다(\Cref{ex:cd-ks-normal}에서 확인). 두 표본,
크기 $n$과 $m$에서는

\[
  1.36\sqrt{\frac{n+m}{nm}}
\]

을 임계값으로 쓴다. \Cref{ex:cd-ks-two}에서 $n=10$, $m=8$이므로

\begin{align*}
  &1.36\sqrt{18/80} = 0.645 && \text{두 표본 근사 임계값(\term{콜모고로프
    극한분포}(Kolmogorov Distribution))}
\end{align*}

$D = 0.6 < 0.645$이므로 이 임계값 기준으로는 기각 실패처럼 보인다.
콜모고로프 극한분포에서 나온 점근 p값 $0.0815$도 같은 결론(기각
실패)이라 이 임계값과 서로 짝이 맞는다. 그러나 표본이 작을 때는
이 점근 근사 자체가 정확하지 않다. 정확 p값을 계산하면 결론이
갈린다.

\begin{pitfall}[슬라이드 값은 "$n \le 6$ 임계값표", 다시 계산하면
"작은 표본에서는 근사와 정확값이 판정을 뒤집을 수 있다"]
슬라이드 23(\texttt{image45})의 임계값표는 한 표본 KS 검정용이고
$n = 1$부터 $6$까지만 보인다. \Cref{ex:cd-ks-normal}($n=100$)이나
\Cref{ex:cd-ks-two}($n=10$, $m=8$)에는 그대로 쓸 수 없다.

두 표본 예제로 이어서 계산하면 다음과 같다. 콜모고로프 극한분포에서
나온 점근 p값은

\[
  P(D \ge 0.6) \approx \texttt{kstwobign.sf}\!\left(0.6\sqrt{80/18}\right)
  = 0.0815
\]

로, $D = 0.6 < 0.645$와 같은 결론(기각 실패)을 준다. 이것이 위
임계값 $0.645$와 짝을 이루는 근사다. 그러나
\texttt{scipy.stats.ks\_2samp(X, Y, method='exact')}로 정확히 계산한
p값은

\[
  p = 2182/43758 = 0.0499
\]

다. $0.0499 < 0.05$이므로 \textbf{기각}이다. 근사(기각 실패)와
정확값(기각)이 판정을 뒤집는다. $0.05$ 경계에서 $0.00013$ 차이로
갈린 결과라, 표본이 작을 때는 근사가 아니라 정확 계산을 써야 한다.
(참고로 \texttt{scipy}의 다른 근사 옵션인 \texttt{method='asymp'}는
$0.0674$를 주는데, 이것은 위 콜모고로프 극한분포가 아니라 유효 표본
크기를 반올림한 또 다른 근사식이며, 결론은 마찬가지로 기각
실패다.)
\end{pitfall}

이 장은 \Cref{ex:cd-ks-two}의 최종 판정을 정확 p값 기준으로
확정한다. $p = 0.0499 < 0.05$이므로 $X$와 $Y$는 같은 분포에서
왔다는 귀무가설을 \textbf{기각}한다.

% =====================================================================
\section{어떻게 적용하는가}
\label{sec:cd-apply}
% source: deck 11 slide 24

\Cref{tab:cd-tests}가 이 장에서 다룬 검정을 한눈에 정리한다.

\begin{table}[ht]\centering
\begin{tabularx}{\linewidth}{@{}l X X X X@{}}
\toprule
비교 대상 & 통계량(Statistic) & 귀무가설(Null Hypothesis) 아래 분포 & 이 장의 예제 & p값 \\
\midrule
평균(2집단) & $t$ & t분포 & \Cref{ex:cd-tomato} & 0.024/0.020 \\
평균(3집단 이상) & $F$ & F분포 & \Cref{ex:cd-cartoon} & 0.008 \\
분포 전체(가설 분포 있음) & $D_n$ & KS 분포(정확) & \Cref{ex:cd-ks-normal} & 0.2997 \\
분포 전체(가설 분포 없음) & $D$ & KS 분포(정확) & \Cref{ex:cd-ks-two} & 0.0499 \\
\bottomrule
\end{tabularx}
\caption{무엇을 비교하느냐에 따라 쓰는 검정 요약. 두 KS 행의 p값은
모두 정확(exact) 계산 값이다. 두 표본 행은 콜모고로프 극한분포에서
나온 점근(asymptotic) 근사 p값 $0.0815$와 결론이 갈린다
(\Cref{sec:cd-ks-pvalue}).}
\label{tab:cd-tests}
\end{table}

\begin{itemize}
  \item \Cref{subsec:ci-compare}: 신뢰구간이 겹칠 때 쓸 "별도의
        검정"이 두 표본 t검정이다(\Cref{sec:cd-t-two}). 버거 대
        치킨너겟 비교는 여기서 새로 계산하지 않는다(치킨너겟의
        표준편차가 본문에 없다).
  \item \Cref{sec:hw2-compare-how}, \Cref{tab:hw2-airport-verdict},
        \Cref{tab:hw2-movie-verdict}: HW2의 $D$ 값은 두 표본 KS
        통계량이다(\Cref{sec:cd-ks-two}). 판정을 $D$ 크기 순위가
        아니라 p값으로 바꾸면 결과가 달라질 수 있지만, 이 장은 HW2
        표에 있는 값만 인용하고 새로 p값을 계산하지 않는다.
  \item \Cref{sec:hw2-learned}: "KS 거리 하나로는 왜 지는지 알 수
        없다"는 슬라이드 23의 요점과 같다. Q-Q 도표와 함께 써야 하는
        이유다.
  \item \Cref{fig:violin}: 보스턴 마라톤 연도별 분포. "여성 주자가
        남성보다 젊은가"(슬라이드 2)를 두 표본 t검정(평균)과
        KS검정(모양)으로 각각 물으면 다른 질문이 된다.
  \item \Cref{ch:project}: 비교 질문을 쓰기 전에 \Cref{fig:cd-choose}로
        어떤 검정을 쓸지 먼저 점검한다.
\end{itemize}
% source: deck 11 slide 24

슬라이드 24는 Kahoot 퀴즈로 끝난다(내용 없음). 대신 확인 문제 셋을
남긴다.

\begin{enumerate}
  \item 질문: "새 비료가 평균 수확량을 늘리는가"에는 어떤 검정이
        맞는가? (답: 두 표본 t검정, \Cref{sec:cd-t-two})
  \item $n = 6$, $t = 2.75$일 때 양측 p값은 어느 범위인가?
        (답: \Cref{ex:cd-t-table}에서 $0.04$ 근방)
  \item $m = 3$, $n = 3$인 ANOVA에서 두 자유도는 무엇인가?
        (답: $\text{df}_{\text{SSB}} = 2$, $\text{df}_{\text{SSW}} = 6$,
        \Cref{ex:cd-cartoon})
\end{enumerate}

% =====================================================================
\section{앞으로}
\label{sec:cd-future}
% source: (이 덱 이후 강의로 이어지는 질문)

\begin{itemize}
  \item \textbf{정규성 가정.} t검정과 F검정은 표본이 정규분포에서
        왔다고 가정한다(슬라이드 3 \texttt{image1}의 "If X is
        normally distributed"). 이 가정은 \term{Q-Q 플롯}(Q-Q Plot)으로
        확인한다.
  \item \textbf{여러 번 검정.} 유전자처럼 수천 번 검정을 반복하면
        기준 $0.05$가 만드는 거짓 양성이 쌓인다(\Cref{sec:pv-threshold}).
        이 덱에는 해법이 없으므로 열린 질문으로 남긴다.
  \item \textbf{모수를 추정한 뒤의 KS 검정과 정규성 전용 검정.} 이
        덱에는 다루지 않으므로 열린 질문으로 남긴다.
  \item \textbf{순열검정}(Permutation Test): 강의계획서가 예고한
        주제(\Cref{sec:ht-future})다. 분포 가정 없이 두 표본을
        비교하는 다른 길로, 이 덱에는 없다.
\end{itemize}
